\section{State Machine and Invariants}
\label{sec:state-machine}
% ================================================================================

\subsection{Product state space}\label{sec:product-model}

The state of a session is a point in the product of orthogonal
dimensions:
\begin{itemize}[leftmargin=2em]
\item \emph{attachment} $a \in \{\textsc{Att}, \textsc{Det},
\textsc{Res}\}$: at least one client is attached, none is, or none is
and a re-attaching client is being verified. It is derived from the
\emph{client set} below;
\item \emph{human} $h \in \{\textsc{None}, \textsc{Pend}\}$: no
escalation is outstanding, or the agent is waiting on a human;
\item \emph{lifecycle} $\ell \in \{\textsc{Active},
\textsc{Terminated}\}$.
\end{itemize}
\paragraph{Client set.} A session has a set $K(s)$ of \emph{registered
client keys}, all bound to the session's \texttt{Owner-ID}: the key
registered at attach, and any key added later by \texttt{CLIENT\_ADDED}.
Its \emph{client set} $C(s) \subseteq K(s)$ is the set of registered keys
that currently have a connected transport (a client is identified by its
key). The attachment is a function of the client set:
$a=\textsc{Att}$ iff $|C| \geq 1$; $a=\textsc{Det}$ iff $|C| = 0$ and
no re-attach is being verified; $a=\textsc{Res}$ iff $|C| = 0$ and a
re-attach request has been accepted and is being verified. A principal
typically has several certified devices, often attached at once;
a second device is a second \emph{join} of the same session, not a
resume. Authority changes (narrowing, I3, I10) apply to the session as a
whole, never to one client. Sharing a session between \emph{different}
principals is out of scope (Section~\ref{sec:scope}).

Every combination of $a$ and $h$ is reachable while the lifecycle is
\textsc{Active}; every terminated session is in the single cell
$\bot$ regardless of $a$ and $h$.

\paragraph{Side predicates.}
The side predicates \texttt{quarantined} and \texttt{key\_exhausted} are
defined over the committed log, alongside the product machine:
\begin{description}[leftmargin=2em]
\item[\texttt{quarantined}$(s)$] (rule~Q): true when the log
contains a \texttt{RESUME\_FAILED} or \texttt{RECOVERY\_FAILED}
event for session $s$, and no subsequent \texttt{TERMINATION}. The
session is frozen; only an operator-issued termination (ending in
\texttt{TERMINATION}) is permitted.
\item[\texttt{key\_exhausted}$(s)$] (I11-REVOKE): true when every
key in $K(s)$ has a \texttt{KEY\_REVOKED} event in the log, and no
subsequent \texttt{OPERATOR\_ADD\_KEY}. No client operation is
accepted; only \texttt{OPERATOR\_ADD\_KEY} (clears the predicate) and
\texttt{OPERATOR\_TERMINATION} (the signed operator form of
\texttt{TERMINATION}, which follows it) are permitted, but the timers keep
running (see Key exhaustion). This predicate is
\emph{not} rule~Q: rule~Q is for storage corruption;
\texttt{key\_exhausted} is for deliberate revocation and has a
recovery path.
\end{description}
Both predicates survive a restart (they are derived from the log,
I12). Neither is a cell of the product machine $Q$.

\begin{definition}[Bridge State Machine]\label{def:sm}
The Bridge Contract state machine is the labelled transition system
$M = (Q, \Sigma, \delta, \lambda, q_0, F)$ where
\begin{itemize}[leftmargin=2em]
\item $Q = (\{\textsc{Att}, \textsc{Det}, \textsc{Res}\} \times
\{\textsc{None}, \textsc{Pend}\}) \cup \{\bot\}$;
\item $\Sigma$ is the set of \emph{inputs} to the bridge: transport
connects and disconnects, re-attach, join and key-addition requests, the
agent's escalation, the human's resolution, the outcome of hash-chain
verification and of recovery, expiry of $T_{\text{detach}}$,
$T_{\text{total}}$ and $T_{\text{human}}$, clean end or abort, and bridge
restart. They are the conditions in the Precondition column of
Table~\ref{tab:transitions}; the Mandatory Governance Events $E$ of
Definition~\ref{def:bc} are the bridge's outputs, not members of $\Sigma$;
\item $\delta : Q \times \Sigma \rightharpoonup Q$ is the partial
transition function defined by Table~\ref{tab:transitions};
\item $\lambda : Q \times \Sigma \rightharpoonup E^{*}$ gives the events
emitted, in order, before the cell changes (the Events column of
Table~\ref{tab:transitions}, I7);
\item $q_0 = (\textsc{Att}, \textsc{None})$;
\item $F = \{\bot\}$.
\end{itemize}
This product machine is the \emph{normative} state model. The
implementation-level notion of a ``state'' used in the rest of this
document is always a cell of $Q$.
\end{definition}

\paragraph{Derived view of named states.} For exposition, and for
compatibility with common usage, the cells are folded into
names: ATTACHED $=(\textsc{Att},\textsc{None})$; DETACHED
$=(\textsc{Det},\textsc{None})$; PENDING\_HUMAN $=(\textsc{Att},
\textsc{Pend})$ or $(\textsc{Det},\textsc{Pend})$; RESUMING
$=(\textsc{Res},\textsc{None})$ or $(\textsc{Res},\textsc{Pend})$;
TERMINATED $=\bot$. Figure~\ref{fig:state-machine} shows this view. The
fold is lossy: it hides whether the client is attached while a human is
pending, which is exactly what determines whether queued tool calls may
drain and whether the detach timeout is running. The derived names are
written in plain capitals; the genesis event of the same spelling is
written \texttt{ATTACHED}. Where the view and
Table~\ref{tab:transitions} disagree, the table is authoritative.

\begin{figure}[t]
\centering
\resizebox{\columnwidth}{!}{% Derived view of named states (Table of cells in the appendix). Drawn from
% the transition table of the product machine; keep the two in step.
\begin{tikzpicture}[
  font=\footnotesize,
  st/.style={draw, rounded corners=3pt, thick, minimum width=2.6cm,
             minimum height=0.9cm, align=center, font=\footnotesize\bfseries},
  ed/.style={-{Stealth[length=2.2mm]}, thick},
  lb/.style={font=\scriptsize, fill=white, inner sep=1pt, align=center},
]
\node[st, fill=blue!8]   (att)  at (0,0)     {ATTACHED};
\node[st, fill=purple!8] (pend) at (8,0)     {PENDING\_HUMAN};
\node[st, fill=gray!20]  (term) at (14,-2.6) {TERMINATED};
\node[st, fill=yellow!20](det)  at (0,-5.2)  {DETACHED};
\node[st, fill=green!10] (res)  at (8,-5.2)  {RESUMING};

\draw[ed] (att) -- node[lb] {last client\\disconnects} (det);
\draw[ed] ([yshift=4pt]att.east) -- node[lb, above=1pt] {escalation} ([yshift=4pt]pend.west);
\draw[ed] ([yshift=-4pt]pend.west) -- node[lb, below=1pt] {human resolves (attached)} ([yshift=-4pt]att.east);
\draw[ed] ([yshift=4pt]det.east) -- node[lb, above=1pt] {re-attach request} ([yshift=4pt]res.west);
\draw[ed] ([yshift=-4pt]res.west) -- node[lb, below=1pt] {client leaves during verification} ([yshift=-4pt]det.east);
\draw[ed] (res.north west) -- node[lb, pos=0.16, xshift=-4pt] {verified,\\no escalation} (att.south east);
\draw[ed, {Stealth[length=2.2mm]}-{Stealth[length=2.2mm]}] (det.north east) -- node[lb, pos=0.16, xshift=4pt] {escalation /\\resolution\\while detached} (pend.south west);
\draw[ed] ([xshift=-4pt]pend.south) -- node[lb, left=1pt] {re-attach request,\\escalation pending} ([xshift=-4pt]res.north);
\draw[ed] ([xshift=4pt]res.north) -- node[lb, right=1pt] {verified,\\escalation pending} ([xshift=4pt]pend.south);
\draw[ed] (pend.east) -- node[lb, pos=0.45, above right] {$T_{\mathrm{human}}$\\exceeded} (term.north west);
\draw[ed] (det.south) -- ++(0,-3.6) -| (term.south);
\node[lb, anchor=south] at (3.3,-9.25) {$T_{\mathrm{detach}}$ or $T_{\mathrm{total}}$ exceeded};

\draw[ed] (att.north west) to[out=120, in=60, looseness=5]
  node[lb, above] {join, add key} (att.north east);
% Quarantine self-loop on RESUMING: drawn on the right side of RESUMING
% so it does not visually overlap the $T_{\mathrm{detach}}$ arrow that
% runs from DETACHED down and across to TERMINATED.
\draw[ed] (res.north east) to[out=60, in=120, looseness=8]
  node[lb, right=2pt, align=left] {verification fails:\\quarantine (rule~Q)} (res.north west);
\end{tikzpicture}
}
\caption{Derived view of named states of the Bridge Contract state machine.
Arrows are allowed transitions; each MUST emit a governance event before
mutating state (Invariant I7). A clean end or abort leads to TERMINATED from
every state, and a successful bridge restart leads to DETACHED; neither is
drawn. A failed verification or recovery quarantines the session and does not
terminate it. The normative model is the product machine of
Definition~\ref{def:sm}.}\label{fig:state-machine}
\end{figure}

In Table~\ref{tab:transitions}, $a$ and $h$ are placeholders for any
value of that dimension unless a value is stated; ``events'' lists the
events emitted, in order, before the cell changes (I7). The table lists
the transitions that change a cell, together with the join and re-attach
transitions that leave it as it is. Events that leave every cell
unchanged (tool-call events, \texttt{KEY\_REVOKED},
\texttt{OPERATOR\_ADD\_KEY}, narrowing, \texttt{AUTHORITY\_EXTENSION\_REJECTED})
are defined by the invariants (I2, I3, I10, I11-REVOKE, I13, I14), not by
a row of the table.

\begin{table}[htbp]
\centering
\caption{Allowed transitions over the product machine.}
\label{tab:transitions}
\footnotesize
\begin{tabularx}{\columnwidth}{l X X}
\toprule
\textbf{Transition} & \textbf{Precondition} & \textbf{Events} \\
\midrule
$(\textsc{Att},h) \to (\textsc{Det},h)$ & Transport disconnect of the
last connected client ($|C|=1$). & \texttt{DETACHMENT},
\texttt{RESUME\_PREPARED}. \\
$(\textsc{Att},h) \to (\textsc{Att},h)$ & Transport disconnect of one
client while others remain ($|C|\geq 2$). & \texttt{CLIENT\_LEFT}. \\
$(\textsc{Det},h) \to (\textsc{Det},h)$ & Idempotent disconnect (also of a
client that is not in $C$). & None. \\
$(\textsc{Att},h) \to (\textsc{Att},h)$ & Join: an authentic request
(I11) from a registered key, with a live unused nonce and a prefix of the
chain (I4, I9); $C := C \cup \{k\}$. & \texttt{CLIENT\_JOINED}; the
events after the acknowledged prefix are replayed to the joiner. \\
$(\textsc{Att},h) \to (\textsc{Att},h)$ & A registered client adds a
further key for the principal. & \texttt{CLIENT\_ADDED}. \\
$(a,\textsc{None}) \to (a,\textsc{Pend})$ & Agent escalates to a
human. & \texttt{HUMAN\_PENDING}; tool queue initialised. \\
$(a,\textsc{Pend}) \to (a,\textsc{Pend})$ & Second escalation while
one is pending: escalations coalesce. & \texttt{HUMAN\_PENDING}, recorded
for audit; the human dimension and $T_{\text{human}}$ (measured from the
first escalation) are unchanged. \\
$(a,\textsc{Pend}) \to (a,\textsc{None})$ & Human resolves. &
\texttt{HUMAN\_RESOLVED}, which ends every coalesced escalation at once.
The queue drains only once $a=\textsc{Att}$ (rule~D below). \\
$(\textsc{Det},h) \to (\textsc{Res},h)$ & Re-attach request that is
signed under a registered client key, uses a live unused nonce, and
names a prefix of the chain (I4, I9, I11). & \texttt{RESUME\_STARTED}. \\
$(a,h) \to (a,h)$, $a \neq \textsc{Res}$ & Re-attach or join request
with a bad signature or unknown key, a consumed, expired or unknown
nonce, a prefix the chain does not contain, or a reported 0-RTT replay
(I4, I8, I9, I11). & \texttt{RESUME\_REJECTED}; no state change. \\
$(\textsc{Att},h) \to (\textsc{Att},h)$ & Re-attach request, authentic
(I11) but made while the session is attached: join is the only way into an
attached session. & \texttt{RESUME\_REJECTED} (\texttt{already\_attached});
no state change, the nonce is spent. \\
$(\textsc{Res},h) \to (\textsc{Res},h)$ & Any further request while one is
being verified. & \texttt{RESUME\_REJECTED} (\texttt{resume\_in\_progress});
no state change. \\
$(\textsc{Res},h) \to (\textsc{Att},h)$ & Hash-chain verification
succeeds. & \texttt{RESUME\_VERIFIED}; then the queue drains (rule~D). \\
$(\textsc{Res},h) \to (\textsc{Det},h)$ & The re-attaching client
disconnects during verification. & \texttt{DETACHMENT}, \texttt{RESUME\_PREPARED}. \\
$(\textsc{Res},h) \to (\textsc{Res},h)$, quarantined & Hash-chain
verification fails (I5). & \texttt{RESUME\_FAILED}; no cell change; the
session is quarantined (rule~Q below). \\
$(\textsc{Det},h) \to \bot$ & $T_{\text{detach}}$ exceeded, or the time
spent detached over all episodes reached $T_{\text{total}}$. &
\texttt{DETACH\_TIMEOUT}, \texttt{TERMINATION}. \\
$(a,\textsc{Pend}) \to \bot$ & $T_{\text{human}}$ exceeded. &
\texttt{HUMAN\_TIMEOUT}, \texttt{TERMINATION}. \\
$(a,h) \to \bot$ & Clean session end or abort. & \texttt{TERMINATION}. \\
$(a,h) \to (\textsc{Det},h)$ & Bridge restart; recovery succeeds
(Section~\ref{sec:recovery}). & \texttt{BRIDGE\_RECOVERED},
\texttt{DETACHMENT}, \texttt{RESUME\_PREPARED}. \\
$(a,h) \to (a,h)$, quarantined & Bridge restart; recovery fails
(Section~\ref{sec:recovery}). & \texttt{RECOVERY\_FAILED}; no cell
change; the session is quarantined (rule~Q below). \\
\bottomrule
\end{tabularx}
\end{table}

\paragraph{Rule D (queue drain).} Queued tool calls drain, in arrival
order, whenever the session is in $(\textsc{Att},\textsc{None})$ with a
non-empty queue: that is, on \texttt{HUMAN\_RESOLVED} while attached, or
on \texttt{RESUME\_VERIFIED} when no escalation is outstanding. A human
who resolves while the client is detached or re-attaching leaves the
queue in place; it drains when the client is back. This is the
conservative reading of G1: the bridge cannot know that a queued call
is client-independent. Each call is re-validated at drain against the current scope, its own
lifetime and approval (I13) and its answer (I14); a call that fails is
not forwarded.

\paragraph{Forwarding acknowledgement.} A drained call passes through
stages that the log records separately: \emph{committed}
(\texttt{TOOL\_CALL\_DRAINED} is durable, I12), \emph{delivered} (the
runtime acknowledges receipt, \texttt{TOOL\_CALL\_DELIVERED}) and
\emph{completed} (the runtime reports the result,
\texttt{TOOL\_CALL\_COMPLETED}). The bridge forwards a call only after
its drain is committed, so a crash can leave a call committed but not
delivered. Such a call stays \emph{unacknowledged}: whenever the session is
in $(\textsc{Att},\textsc{None})$ and a bridge instance has not yet
forwarded an unacknowledged call, it forwards the call again with the
same idempotency key (I13), without a second drain event. A call that is
delivered is never forwarded again. Delivery is therefore at-least-once
with deduplication by the runtime; exactly-once execution is the runtime's
responsibility.

\paragraph{Calls outside an escalation (G1) --- conformance profiles.}
Rule D and I2 govern calls
that arrive while $h=\textsc{Pend}$. A call that arrives with
$h=\textsc{None}$ while $a\neq\textsc{Att}$ has no escalation to wait for,
and the contract gives the deployment a choice among the
\emph{conformance profiles} \texttt{refuse}, \texttt{queue} and
\texttt{forward}, fixed by configuration
and the same for every session of the deployment. The deployment
documents which profile it uses and, for the \texttt{forward} profile,
who may configure the forward allowlist (an operator role named in the
deployment documentation).
\begin{description}[leftmargin=2em]
\item[\texttt{refuse} (mandatory default).] The bridge refuses the call: it is not
queued, not forwarded and not recorded as a queued call, and the error
reaches the runtime. This profile is mandatory: every conformant
bridge MUST implement it, and it is the default if the deployment does
not name one. Every other statement in
this section assumes this profile unless it says otherwise.
\item[\texttt{queue} (optional).] The bridge queues the call
(\texttt{TOOL\_CALL\_QUEUED} with \texttt{held} set). It is held until a
human answers it (I14) and the session is in $(\textsc{Att},\textsc{None})$,
and is then forwarded under rule~D. A held call that is still unanswered
at drain is not rejected as unanswered; it waits and expires after
$T_{\text{call}}$ (I13). Held calls do not depend on an escalation, so
$T_{\text{approve}}$ does not apply to them.
\item[\texttt{forward} (optional).] The operator lists permissions that may run
without the client. A call whose permission is on that list and in the
current scope is recorded as queued, answered (source
\texttt{operator\_policy}) and drained in one committed batch, and
forwarded after the batch is durable (I12). Any other call is refused as
under \texttt{refuse}.
\end{description}
No profile changes I2: while $h=\textsc{Pend}$ every call is queued or
rejected under rule~D, whatever the profile. The \texttt{forward} profile
is the operator-defined exception to rule~D that a deployment opts into;
the bridge cannot tell on its own which calls are client-independent.

\paragraph{Rule Q (quarantine).} A session whose log cannot be vouched
for is \emph{quarantined}, not terminated. This applies when the hash
chain fails to verify after a re-attach was accepted
(\texttt{RESUME\_FAILED}), and when recovery after a restart finds the
chain broken
(\texttt{RECOVERY\_FAILED}); the cause is the storage or an attacker on it,
not the client, so ending the session would punish the wrong party and
destroy the evidence. The bridge MUST record the failure if it can still
append, MUST NOT emit \texttt{TERMINATION}, and MUST then refuse every
client operation on the session (attach, join, challenge, escalation,
tool call, authority change), run no timer on it and emit no further event
for it, except an operator's explicit termination (\texttt{TERMINATION}).
An operator's key addition is refused on a quarantined session
(\texttt{quarantined}), and no \texttt{CHECKPOINT} is appended to one. The session's
event log MUST be preserved read-only for forensic review: the bridge
MUST NOT append, truncate, or rewrite any record of a quarantined
session except the failure event that rule~Q itself permits and any
later \texttt{TERMINATION} the operator issues. The session
keeps the cell it had; quarantine is not a cell of the product machine,
and a restart does not lift it (the failure event is in the log; if it
could not be appended, recovery finds the broken chain again). Release
from quarantine is an operator action outside the contract; the only
action the contract permits on a quarantined session is operator
\texttt{TERMINATION}. Chain
verification follows authentication (I11), so a presenter without a
registered key cannot cause a quarantine. Quarantine is defence in depth
for a violated A1 (it freezes a session whose chain no longer verifies);
the guarantee of the safety argument rests on A1 and does not include it.

\subsection{Conformance Invariants}

\begin{invariant}[I1 --- Resume Preparedness]\label{inv:I1}
For every transition into a cell with $a=\textsc{Det}$ caused by a
disconnect (of the last connected client from $(\textsc{Att},h)$, or of
the re-attaching client from $(\textsc{Res},h)$), the bridge
MUST emit \texttt{DETACHMENT} followed by \texttt{RESUME\_PREPARED}
\emph{before} the cell changes.
\end{invariant}

\paragraph{Termination effects.}
On \texttt{TERMINATION}, the session enters $\bot$ (I6). The
following effects are normative:
\begin{itemize}[leftmargin=2em]
\item The tool-call queue is discarded; queued calls are not
forwarded. The runtime MUST be notified that the session is
terminated and SHOULD clean up any in-flight work.
\item In-flight tool calls (calls the bridge forwarded before
termination but whose \texttt{TOOL\_CALL\_COMPLETED} was not
emitted) are not tracked further; the runtime is responsible for
their disposition.
\item No timer runs on a terminated session; no further events are
emitted; no client operation is accepted (I6).
\item The event log is retained for audit (A1); the deployment
SHOULD apply its retention policy (Section~\ref{sec:limitations}).
\end{itemize}

\begin{invariant}[I2 --- Tool Call Suspension]\label{inv:I2}
While $h=\textsc{Pend}$, every tool call that the runtime submits to the
bridge for approval or forwarding MUST be either queued (with \texttt{TOOL\_CALL\_QUEUED} emitted) or rejected.
The bridge MUST NOT forward any tool call to the agent runtime for
execution while $h=\textsc{Pend}$. In-flight tool calls initiated
before \texttt{HUMAN\_PENDING} are not retracted; their completion is
audited as \texttt{TOOL\_CALL\_COMPLETED} (the separate completion event
of the forwarding acknowledgement, rule~D).
I2 covers what the runtime submits to the bridge; execution that the
bridge does not mediate is outside the contract (Section~\ref{sec:coverage}).
\end{invariant}

\begin{invariant}[I3 --- Authority Lock]\label{inv:I3}
In every cell, any attempt to widen \texttt{Authority-Scope} (any
proposed scope containing a permission that the current scope lacks)
MUST be rejected and MUST emit \texttt{AUTHORITY\_EXTENSION\_REJECTED};
the scope the bridge holds is unchanged. The scope of a session may only
narrow (I10). Widening requires a new attach with a freshly granted
scope, which is a new session. A tool call submitted while the session is
in $(\textsc{Att}, \textsc{None})$ with no human decision pending, whose
permission is not in the current scope, MUST be refused with the reason
\texttt{scope} and recorded as \texttt{TOOL\_CALL\_REJECTED} with
\texttt{direct} set; a call inside the scope is forwarded as before. I3 has two parts:
\begin{description}[leftmargin=2em]
\item[Enforced.] A change requested \emph{through the bridge} is refused
before it takes effect.
\item[Observed.] A change that the runtime announces \emph{itself} (for
example an Agent Client Protocol (ACP) agent switching its own mode; see
Section~\ref{sec:rw-agtp-thread} for the ACP disambiguation) has already happened when
the bridge learns of it. The bridge MUST emit
\texttt{AUTHORITY\_EXTENSION\_REJECTED} with source \texttt{observed}
and keep its record unchanged, and SHOULD instruct the runtime to revert
to the recorded scope. I3 does not claim that the bridge prevents a
runtime-initiated widening, only that it never records one as granted and
always reports it.
\end{description}
\end{invariant}

\begin{invariant}[I4 --- Prefix consistency]\label{inv:I4}
A re-attach or join request carries the pair $(\mathit{seq}, \mathit{hash})$ of
the last event the client acknowledged, taken from the event stream the
client observed, not from any artefact issued at detach. The bridge
MUST accept the request only if $\mathit{seq}$ does not exceed the
sequence number of the last committed event and $\mathit{hash}$ equals
the hash of the event with sequence number $\mathit{seq}$ in the bridge's
chain: that is, the client's history is a \emph{prefix} of the bridge's
chain. Any other pair (a different hash at that sequence, or a sequence
the chain has not reached) MUST be rejected with
\texttt{RESUME\_REJECTED} (reason \texttt{prefix\_mismatch}). Rejection
leaves the session in its current cell; it does not terminate the
session. On acceptance the bridge MUST deliver every committed event
after $\mathit{seq}$ to the client (the re-attaching or joining client),
including those that changed the
authority scope or the human dimension while the client was away; the
client reconciles its view from this replay.
\end{invariant}

\begin{invariant}[I5 --- Hash Chain Integrity]\label{inv:I5}
Every Governance Event MUST carry a \texttt{prev\_hash} field equal
to the SHA-256 of the canonical JSON of the immediately preceding
event in the same session. The genesis event's \texttt{prev\_hash}
MUST be the empty string.

I5 detects an insertion, deletion, rewrite or reordering of events in
the middle of the log. It does not detect the removal or rewrite of the
last events: a log cut back to an earlier event, or whose final event is
replaced, is again a valid chain. Detecting that requires the head
(sequence and hash) to be committed outside the log, as an
I-CHECKPOINT anchor does.
\end{invariant}

\begin{invariant}[I6 --- Terminality]\label{inv:I6}
The cell $\bot$ is terminal. Any subsequent operation that would
mutate state MUST raise an error. This covers every operation, including
attach, disconnect, join, add a client key, escalate, re-attach, submit
tool call, attempt authority extension, narrowing, human resolution, a
tool-call answer and the operator's key operations.
\end{invariant}

\begin{invariant}[I7 --- Event-before-Mutation]\label{inv:I7}
For every transition, the corresponding governance events MUST be
emitted and appended to the event log \emph{before} the cell is
changed. This ensures that the audit trail records the intent to
transition before the transition takes effect. Together with I12 this
ordering survives a bridge crash.
\end{invariant}

\subsection{Assumptions and Axioms}\label{sec:axioms}

The safety argument in Section~\ref{sec:safety-argument} rests on the
following assumptions, stated as axioms. A bridge implementation
operating in violation of any axiom is out of scope of this
specification.

\begin{description}[leftmargin=2em]
\item[A1 --- Reliable storage.] A record that has been committed under
I12 is not lost, reordered, or mutated by the storage medium or by the
bridge itself. (Mid-log tampering is detected by I5; truncation or
rewriting of the last events is not, without an external head.)
Rule~Q is defence in depth for a violated A1 and is not part of the
guarantee that rests on A1.
\item[A2 --- Monotonic sequence by construction.] The bridge's
sequence-number generator returns a strictly increasing value
on each call, including across restarts (it is recovered from the
log). Production of sequence numbers is monotonic by construction;
acceptance of client-claimed sequences is governed by I4.
\item[A3 --- SHA-256 collision resistance.] SHA-256 is collision
resistant under the assumed adversary. \texttt{prev\_hash} chaining
and the prefix check of I4 rely on this. The signature scheme of I11 is
assumed to be existentially unforgeable, and nonces are drawn from a
cryptographically secure source.
\item[A4 --- Trusted bridge.] The bridge is non-Byzantine: it does not
deliberately suppress governance events, forge events,
or lie about its state. (Byzantine bridge mitigations
are deployment-level, see \S\ref{sec:scope}.)
\item[A5 --- Serialised event emission per session.] Concurrent
event emission per session is serialised by the bridge's internal
locking. Concurrent sessions are independent. Because the dimensions of
the product machine are orthogonal, serialised events from different
dimensions (e.g., \texttt{DETACHMENT} and \texttt{HUMAN\_PENDING})
commute on the cell they reach.
\item[A6 --- Clocks.] While the bridge runs, timeouts use a monotonic
clock; wall-clock skew does not affect them. Across a restart, elapsed
time is taken from the timestamps of committed events, so a timeout may
fire late by at most the wall-clock error of the host. A host clock that
reads earlier than the last committed timestamp is treated as reading
that timestamp: no time is counted until the clock passes it, and
committed timestamps never decrease. Downtime counts toward
$T_{\text{human}}$, which is measured from the original
\texttt{HUMAN\_PENDING}, and toward $T_{\text{total}}$, but not toward
$T_{\text{detach}}$, which restarts at the recovered detachment
(Section~\ref{sec:recovery}). Fair scheduling
is assumed: the timeout check is eventually invoked (see
Appendix~\ref{app:formal-limitations}).
\item[A7 --- Key custody.] The private keys of registered clients are
not disclosed to the adversary, and the registry of client public keys
survives restarts (it is rebuilt from the log, I12). Compromise of a
client key is a deployment-level concern. The contract provides
operator-signed revocation (I11-REVOKE): a lost or compromised device
key can be revoked by the operator, after which it is no longer
accepted for re-attach or join. Soft revocation (\texttt{rotation})
keeps active connections; hard revocation (\texttt{compromise})
removes the key from the client set and detaches the session if that was
the last connected client. When all keys are revoked, the session enters
\texttt{key\_exhausted} (not rule~Q); the operator can add a new key
or terminate, and otherwise a detached session ends after the detachment
bound. Before revocation, a lost device stays authorised
until the session ends, which in the persistent profile can be
days; revocation is the operator's recovery path.
\end{description}

\subsection{Deployment requirements}
\label{sec:deployment-requirements}

The axioms above state what the contract assumes of the bridge and the
storage. A conformant deployment must additionally provide what
Table~\ref{tab:deployment-requirements} lists; the table makes explicit
what the baseline MUSTs leave to the operator.

\begin{table}[htbp]
\centering
\caption{Assumptions and deployment requirements.}
\label{tab:deployment-requirements}
\footnotesize
\begin{tabularx}{\columnwidth}{l X}
\toprule
\textbf{Requirement} & \textbf{What the deployment must provide} \\
\midrule
Append-only storage &
A storage medium that, once a record is committed under I12, does not
lose, reorder, or mutate it (A1). Mid-log tampering is detected by I5;
truncation of the tail is detected only against a head held outside the
log, and, if required, by I-SIGN-EVENT and I-CHECKPOINT. \\
Trusted bridge &
A non-Byzantine bridge (A4): one that does not deliberately suppress
governance events, forge events, or lie about its state. Byzantine-bridge
mitigations are deployment-level (external transparency logs). \\
Client-key custody &
The principal holds the private keys of every registered client key
(A7); the contract provides operator-signed revocation (I11-REVOKE)
for lost or compromised keys. Soft revocation (\texttt{rotation})
keeps active connections; hard (\texttt{compromise}) removes the key
from the client set and detaches the session if it was the last
connected client. When
all keys are revoked, the session enters \texttt{key\_exhausted};
the operator adds a new key or terminates, and otherwise a detached
session ends after the detachment bound. \\
Human-resolution and answer channel &
An authenticated channel for \texttt{HUMAN\_RESOLVED},
\texttt{TOOL\_CALL\_ANSWERED} and a clean end of session. The deployment
documents who may emit them and under which key. The contract does not
specify the channel; the bridge trusts whatever the deployment feeds it. \\
Runtime idempotency contract &
The runtime deduplicates forwarded calls by idempotency key. The key
scope is the session. The deployment states a retention bound for
spent keys. A second submission of the same key with different content
is rejected: this is a MUST while the first call is queued (I13,
reason \texttt{duplicate\_key}); for a key already spent, the same
rejection is a SHOULD that holds within the deployment's retention
bound, and beyond that bound the contract states no requirement.
\texttt{DELIVERED}
means receipt, not execution; exactly-once external effects are not
guaranteed. \\
Runtime routes execution and authority through the bridge &
Where the deployment wants the bridge to enforce (rather than observe)
authority, the runtime routes every authority change and every
tool-call submission through the bridge. Execution the bridge does not
mediate is outside the contract (Section~\ref{sec:coverage}). \\
Log access authorisation and retention &
The deployment documents who may read the event log and the replay, and
the bridge enforces a retention bound after which the log is archived or
destroyed (Section~\ref{sec:limitations}). \\
\bottomrule
\end{tabularx}
\end{table}

\subsection{Lemmas}\label{sec:lemmas}

The following lemmas are immediate consequences of the invariants and
axioms. Their proofs are short and are provided for completeness.

\begin{lemma}[Audit ordering]\label{lem:atomicity}
I5 $+$ I7 $+$ I12 $\Rightarrow$ no externally observable change of cell
exists without its governing events already committed in the log, and
the log alone determines the cell that recovery rebuilds.
\end{lemma}

\begin{proof}[Sketch.]
I7 requires the events to be appended before the cell changes, and I12
requires the append to be durable before the change and before any
acknowledgement to the client. I5 links each event to its predecessor.
Therefore any cell visible to a client, or recoverable after a crash,
is preceded by a verifiable committed event. A crash between the append
and the change leaves the log ahead of the cell; recovery
(Section~\ref{sec:recovery}) resolves this in favour of the log. The
lemma asserts ordering, not atomicity: it does not claim that append and
change occur indivisibly. \qedhere
\end{proof}

\begin{lemma}[No authority drift]\label{lem:no-drift}
I3 $+$ I10 $\Rightarrow$ in every cell, at every point of an execution,
the authority scope that the bridge holds for a session is a subset of
its scope at attach. (The runtime's effective authority equals it only
for changes requested through the bridge; see I3, observed.)
\end{lemma}

\begin{proof}[Sketch.]
The scope is changed only by an accepted narrowing (I10); every
widening attempt is rejected without changing it (I3). By induction
over the transitions, the scope after each step is a subset of the scope
before it, hence of the scope at attach. \qedhere
\end{proof}

\begin{lemma}[Resume recoverability]\label{lem:recover}
I1 $+$ I4 $+$ I5 $+$ I11 $\Rightarrow$ on an authenticated re-attach,
the client is brought to the bridge's current state by a verifiable
replay of the events it has not seen.
\end{lemma}

\begin{proof}[Sketch.]
I1 forces \texttt{RESUME\_PREPARED} to be emitted and stored before the
cell changes to $a=\textsc{Det}$, so the resume point is recorded
durably. I11 ensures the presenter holds a registered client key and
answers a fresh nonce, so the request is neither forged nor replayed. I4
ensures that the $(\mathit{seq}, \mathit{hash})$ the client presents lies
on the bridge's chain, and I5 ensures that every event after it carries
the \texttt{prev\_hash} of its predecessor, so the client can verify the
replay against the hash it already holds. Hence the client's resulting
view is a causally consistent extension of what it had acknowledged.
\qedhere
\end{proof}

\begin{lemma}[Terminality]\label{lem:terminality}
I6 $\Rightarrow$ after $\bot$, no further state mutation is
possible for the session.
\end{lemma}

\begin{proof}[Sketch.]
The lemma restates I6: I6 forbids every state-mutating operation on a
session in $\bot$, and $\bot$ has no outgoing transition in
Table~\ref{tab:transitions} ($F=\{\bot\}$). \qedhere
\end{proof}

\subsection{Liveness properties}\label{sec:liveness}

The Bridge Contract's liveness properties depend on the timeouts
defined in \S\ref{sec:timeouts}. Without timeouts the contract
satisfies safety but not liveness: a session could remain in a cell
with $a=\textsc{Det}$ or $h=\textsc{Pend}$ indefinitely. The following
liveness properties hold \emph{conditionally on the bridge's timer
mechanism being scheduled fairly} (Axiom A6; see
Appendix~\ref{app:formal-limitations} for the fairness limitation in the
formal model). Without a fairness assumption on the scheduler, L1 and L2
do not hold: a scheduler that indefinitely postpones the timeout check
would leave the session in such a cell forever, which does not violate
any safety invariant but violates liveness.

\begin{description}[leftmargin=2em]
\item[L1.] \emph{(conditional on fair scheduling, A6)} For every
session that enters a cell with $a=\textsc{Det}$, either the client
re-attaches successfully (through $a=\textsc{Res}$ to
$a=\textsc{Att}$), or the session reaches $\bot$ within
$T_{\text{detach}}$ of that detachment, or it is quarantined (the exception
discussed after L2). A \texttt{key\_exhausted} session is no exception: its
timers keep running, so unless the operator adds a key it reaches $\bot$
within $T_{\text{detach}}$ as well. Over the whole session, the time
spent in cells with $a=\textsc{Det}$, summed over all detachments and
including time during which the bridge was down, never exceeds
$T_{\text{total}}$: once it has used up $T_{\text{total}}$, a session in
$a=\textsc{Det}$ reaches $\bot$ even if no single detachment has lasted
$T_{\text{detach}}$ (the model checks that it then leaves $a=\textsc{Det}$,
by a re-attach in progress, or ends, or is quarantined; \texttt{CapEnds}).
A re-attach does not depend on any artefact issued
at detach (I4, I11), so no secondary expiry can terminate the session
earlier.
\item[L2.] \emph{(conditional on fair scheduling, A6)} For every session
that enters a cell with $h=\textsc{Pend}$, either the human resolves
(the session moves to $h=\textsc{None}$) or the session reaches
$\bot$ within $T_{\text{human}}$ of the first escalation, or earlier if
the detach timeout fires first, or the session is quarantined. The
human timer of a \texttt{key\_exhausted} session keeps running. While $a=\textsc{Det}$ the detach
timer runs (it is stopped in $a=\textsc{Res}$), so a human wait never outlasts the detachment the deployment
tolerates; while the client is attached the wait may be longer than
$T_{\text{detach}}$ if $T_{\text{human}}$ is.
\end{description}

L1 and L2 together imply that no active session that is not quarantined
remains in a transient cell indefinitely (under fair scheduling), with one
reservation about a \texttt{key\_exhausted} session: while a client is
still connected on a softly revoked key the cell is $a=\textsc{Att}$ and no
detach timer runs, until the transport is lost or the operator acts (the
hazard of rotating the only key is described under Key exhaustion).
$\bot$ is the only
cell in which a session may persist without bound, apart from
$(\textsc{Att},\textsc{None})$, which has no timer and is stable by design,
and from a quarantined session, which is frozen by rule~Q until an operator
acts. A re-attach
cannot hold the session in $a=\textsc{Res}$ without bound, because a
client disconnect during verification returns it to $a=\textsc{Det}$ and
verification itself is a finite computation over the log. A re-attach that
is started and abandoned repeatedly is a limit of this bound: time in
$a=\textsc{Res}$ and $a=\textsc{Att}$ is not counted by $T_{\text{total}}$,
so $T_{\text{total}}$ does not bound such a cycle. A revocation of the only key while it is being
verified, soft or hard, abandons the verification and the session returns
to $a=\textsc{Det}$, where the detach timer runs.

\subsection{Timeouts}\label{sec:timeouts}

The state machine depends on the timers of Table~\ref{tab:timers}.
$T_{\text{human}}$, $T_{\text{call}}$ and $T_{\text{total}}$ default from
$T_{\text{detach}}$; $T_{\text{approve}}$ and $T_{\text{nonce}}$ do not
(Table~\ref{tab:timer-profiles}). The
applicability of $T_{\text{detach}}$ and $T_{\text{human}}$ to each cell
is:

\begin{center}
\small
\begin{tabular}{l l l}
\toprule
\textbf{Cell} & $T_{\text{detach}}$ & $T_{\text{human}}$ \\
\midrule
$(\textsc{Att},\textsc{None})$ & --- & --- \\
$(\textsc{Att},\textsc{Pend})$ & --- & running \\
$(\textsc{Det},\textsc{None})$ & running & --- \\
$(\textsc{Det},\textsc{Pend})$ & running & running \\
$(\textsc{Res},\textsc{None})$ & stopped & --- \\
$(\textsc{Res},\textsc{Pend})$ & stopped & running \\
\bottomrule
\end{tabular}
\end{center}

\begin{table}[htbp]
\centering
\caption{The timers: start event, clock source, pause/reset
behaviour, survival across a bridge restart, expiry event and
interactions. Defaults and bounds are listed in
Table~\ref{tab:timer-profiles}.}
\label{tab:timers}
\scriptsize
\begin{tabularx}{\columnwidth}{l X X X X X}
\toprule
\textbf{Timer} & \textbf{Starts on} & \textbf{Clock} & \textbf{Pause / reset} & \textbf{Survives restart?} & \textbf{Expiry event} \\
\midrule
$T_{\text{detach}}$ &
The \texttt{DETACHMENT} that begins a detachment; not restarted by later events while detached. &
Monotonic while running; across a restart, elapsed time is taken from committed event timestamps (A6). &
Stopped while $a=\textsc{Res}$; a return from $a=\textsc{Res}$ to $a=\textsc{Det}$ is a new \texttt{DETACHMENT} and starts a fresh interval ($T_{\text{total}}$ keeps counting the time spent detached, so that time is bounded over the cycle; a re-attach started and abandoned repeatedly is not bounded by $T_{\text{total}}$, which does not count time in $a=\textsc{Res}$); reset to a fresh interval by a bridge restart while detached (recovery rule~2). &
The interval is reset (the restart begins a new detachment). &
\texttt{DETACH\_TIMEOUT}, \texttt{TERMINATION}. \\
$T_{\text{total}}$ &
The first \texttt{DETACHMENT} of the session; summed over every away interval. &
Computed from committed event timestamps; never held as an in-memory timer. &
Never reset; time in $a=\textsc{Res}$ and $a=\textsc{Att}$ is not counted; bridge downtime while detached is counted. &
Yes --- downtime is part of the away interval. &
\texttt{DETACH\_TIMEOUT}, \texttt{TERMINATION}. \\
$T_{\text{human}}$ &
The first \texttt{HUMAN\_PENDING} of the current escalation; a second escalation while one is pending does not restart it. &
Monotonic while running; across a restart, elapsed time is taken from the original \texttt{HUMAN\_PENDING} (A6). &
Running while $h=\textsc{Pend}$; stopped when $h=\textsc{None}$; bridge downtime while $h=\textsc{Pend}$ counts toward it. &
Yes --- measured from the original \texttt{HUMAN\_PENDING}; an elapsed timeout is applied after recovery. &
\texttt{HUMAN\_TIMEOUT}, then \texttt{TERMINATION}. \\
$T_{\text{call}}$ &
The \texttt{TOOL\_CALL\_QUEUED} event of the call. &
Monotonic. &
Not paused; expires even if the queue is held under the \texttt{queue} profile. The call is removed on expiry. &
Yes (queued in the log); a call whose $T_{\text{call}}$ elapsed during downtime is expired at recovery. &
\texttt{TOOL\_CALL\_EXPIRED}; the call is not forwarded. \\
$T_{\text{approve}}$ &
The \texttt{HUMAN\_RESOLVED} event of the escalation. &
Monotonic. &
Not paused; runs even while the client is away. An approval that lapses is void at the next drain. &
Yes (queued in the log). &
\texttt{TOOL\_CALL\_EXPIRED} on every call of the lapsed escalation at the next drain; the expiry is recorded per call when the approval is given, so a later approval does not revive it. \\
$T_{\text{nonce}}$ &
The bridge issues a challenge and returns a nonce. &
Monotonic; the bridge's own clock. &
Not paused; the nonce is volatile and is void after a restart. &
No --- outstanding nonces are void; the client asks for a fresh challenge. &
The nonce is rejected on next use with \texttt{RESUME\_REJECTED} (reason \texttt{nonce\_expired}). \\
\bottomrule
\end{tabularx}
\end{table}

\paragraph{Derivation of the timers.} The timers are not independent:
\[
T_{\text{call}} \;\leq\; T_{\text{human}}
\qquad\text{and}\qquad
T_{\text{detach}} \;\leq\; T_{\text{total}}
\qquad (\text{MUST}),
\]
and $T_{\text{human}}$ defaults to $T_{\text{detach}}$,
$T_{\text{call}}$ to $T_{\text{human}}$ and $T_{\text{total}}$ to four
times $T_{\text{detach}}$. $T_{\text{call}} \leq T_{\text{human}}$ is a configuration simplification:
a call may not outlive the longest wait the session tolerates (under the
\texttt{queue} profile a held call has no escalation, but is bounded the
same way). The cap is useless if it is below a single detachment.
$T_{\text{human}} \leq T_{\text{detach}}$ is a recommendation that the
profiles follow, not a configuration bound: it constrains a session only
where the client is away, and there the detach timer enforces it without
help, because a detached session ends at $T_{\text{detach}}$ whatever the
human timer says. With the client attached, an escalation may wait longer
than $T_{\text{detach}}$ if the deployment sets $T_{\text{human}}$ so. When
the client then detaches, the wait is bounded by the time that remains of
$T_{\text{human}}$ and by $T_{\text{detach}}$ from the detachment,
whichever ends first. Independent defaults (for example 24~h and 1~h) that
did not follow from $T_{\text{detach}}$ would end work that a mobile owner
expects to survive a long detach. A deployment chooses one
\emph{profile}:

\begin{table}[htbp]
\centering
\caption{Timer profiles. Defaults are informative; operators MAY set
the timers to any non-negative values that satisfy the bounds above (the
value 0 is degenerate: a nonce expires on issue, a detachment ends at
once).}
\label{tab:timer-profiles}
\footnotesize
\begin{tabular}{l r r r r r r}
\toprule
\textbf{Profile} & \makebox[1.4cm]{$T_{\text{detach}}$} & \makebox[1.4cm]{$T_{\text{human}}$} &
\makebox[1.2cm]{$T_{\text{call}}$} & \makebox[1.4cm]{$T_{\text{total}}$} & \makebox[1.3cm]{$T_{\text{approve}}$} &
\makebox[1.1cm]{$T_{\text{nonce}}$} \\
\midrule
interactive-short (default) & 30 min & 30 min & 30 min & 2 h & 15 min & 60 s \\
persistent & 7 days & 7 days & 7 days & 28 days & 15 min & 60 s \\
\bottomrule
\end{tabular}
\end{table}

\noindent $T_{\text{approve}}$ (I13) is a freshness bound on an approval,
not a liveness bound on the session, and is deliberately independent: in
the persistent profile an approval given while the client is away still
expires after 15 minutes and the call is rejected at drain, so the agent
must ask again. A deployment that wants approvals to outlive a detachment
raises $T_{\text{approve}}$ and accepts the weaker freshness.

\paragraph{Interaction.} In $(\textsc{Det},\textsc{Pend})$ both timers
run and whichever fires first terminates the session; the cumulative cap
$T_{\text{total}}$ runs in the same cells and ends the session if it is
reached first. Because a detached session ends at $T_{\text{detach}}$
regardless, holding it open for a longer $T_{\text{human}}$ would serve
no purpose. A deployment that must wait on a human for longer while the
client is away MUST raise $T_{\text{detach}}$ (and $T_{\text{total}}$)
accordingly.

\paragraph{Conflict resolution (both timeouts fire simultaneously).}
If both timeouts expire in the same tick in
$(\textsc{Det},\textsc{Pend})$, the bridge MUST:
\begin{enumerate}[leftmargin=2em]
\item Emit \texttt{HUMAN\_TIMEOUT} (the human-side timeout event);
\item Emit \texttt{DETACH\_TIMEOUT} (the detach-side timeout event,
to preserve audit completeness --- suppressing it would lose
information about which timeouts fired, contradicting the
event-completeness rationale of Table~\ref{tab:event-completeness});
\item Emit \texttt{TERMINATION} with reason field
\texttt{human\_timeout+detach\_timeout (simultaneous)};
\item Transition to $\bot$.
\end{enumerate}
The order is fixed: \texttt{HUMAN\_TIMEOUT}, then
\texttt{DETACH\_TIMEOUT}, then \texttt{TERMINATION}, so that two logs of the
same run are equal. The choice of \texttt{human\_timeout} as the leading
reason-token in (3) follows a default priority rule, not an
alphabetical one: the human escalation is the semantically more salient event for
audit consumers. The bridge
SHOULD apply this default consistently across all simultaneous-timeout
cases; operators MAY override the default priority via configuration
(e.g., to prefer \texttt{detach\_timeout} in deployments where
detach is the more salient signal). An override changes only the leading
token of the reason in (3), never the order of the events. An override is a
deployment configuration, not a contract violation.

\paragraph{Edge cases (normative).}
\begin{enumerate}[leftmargin=2em]
\item \emph{Timer expiry in the same tick as an incoming request.}
If a timer expires in the same tick as an incoming request that would
move the session out of the timed cell (for example, $T_{\text{detach}}$
expiring in the same tick as a re-attach request that names a valid
prefix), the bridge MUST process the request first. If the request
succeeds (e.g., \texttt{RESUME\_STARTED}), the timer is no longer
applicable to the new cell and is not fired. If the request fails
(e.g., \texttt{RESUME\_REJECTED}), the timer fires in the same tick
and the session transitions accordingly. The reference bridges follow the
rule: a timer that expires in the tick of a re-attach or a join waits for the
request, and one that expired in an earlier tick has ended the session, so
the request is refused uniformly. The one exception is $T_{\text{human}}$: it
wins over a resolve that arrives in its own tick (a resolve is the only request
that moves a session out of the pending-approval cell; every other request
leaves the session in that cell, so the rule above does not apply to it). The bridge ends the session
with \texttt{HUMAN\_TIMEOUT} and \texttt{TERMINATION}, refuses the resolve and
answers it \texttt{terminated}.
\item \emph{The same idempotency key arriving with different content.}
A second tool call submitted with an idempotency key that is already
queued MUST NOT be queued, regardless of whether the content matches
or differs from the first submission; the caller is answered with the
reason \texttt{duplicate\_key}. The first submission's content
and hash are retained; the second is discarded and commits no event, so
the queued record is returned and the log does not record the second
submission (limitations register).
\end{enumerate}

\paragraph{Configuration.} Default values are informative.
Operators MAY set the timers to any values $\geq 0$ that satisfy the
bounds above. A timeout of $0$ means ``immediate termination on entry to the cell''; this is
useful for testing and for deployments that prefer fail-fast.
With $T_{\text{detach}} = 0$: on entry to a cell with $a=\textsc{Det}$,
the bridge MUST still emit \texttt{DETACHMENT} +
\texttt{RESUME\_PREPARED} (per I1) before transitioning to $\bot$ with
\texttt{DETACH\_TIMEOUT} + \texttt{TERMINATION}.

\subsection{Bridge restart and recovery}\label{sec:recovery}

When the bridge process restarts with live sessions, it MUST recover
each session from its durable event log (I12) before accepting any
client message for it. The rules are:

\begin{enumerate}[leftmargin=2em]
\item \textbf{Replay.} Read the committed log from genesis (or from the
latest verified checkpoint, I-CHECKPOINT), verify the hash chain (I5),
and rebuild the cell, the scope and \texttt{scope\_version}, the
registry of client keys, the tool queue, the idempotency keys, and the
timer anchors. A trailing record
that is not fully durable is discarded; by I12 it has no externally
visible effect.
\item \textbf{Transports are gone.} Every client transport was lost with
the process, so every recovered active session has an empty client set
and $a=\textsc{Det}$; the registered keys $K$ are rebuilt from the log,
with the scheme tag and public key each event records, and the key registry
in storage is repaired from them: the log is the source of truth, so a key
missing from the registry is written back, an entry the log does not record
(a crash between the registry write and the event commit leaves one) is
dropped, and an entry whose scheme or public key differs from the log is
replaced. \texttt{BRIDGE\_RECOVERED} names the key ids repaired.
The bridge MUST emit \texttt{BRIDGE\_RECOVERED}, then \texttt{DETACHMENT}
(reason \texttt{bridge\_restart}) and \texttt{RESUME\_PREPARED} (the
resume record of the new detachment), before the cell changes.
$T_{\text{detach}}$ restarts at the new \texttt{DETACHMENT}; the
cumulative cap $T_{\text{total}}$ does not (a restart while detached
continues the away interval, downtime included). Challenge
nonces are volatile and are void after a restart: a client that holds
one simply requests another (I11).
\item \textbf{Elapsed timeouts.} $T_{\text{human}}$ is measured from the
original \texttt{HUMAN\_PENDING} (A6). If it has already elapsed, the
bridge MUST terminate the session after recovery with
\texttt{HUMAN\_TIMEOUT} and \texttt{TERMINATION}. Likewise, if the
cumulative cap $T_{\text{total}}$ has already been reached (downtime
included), the reference bridge terminates the session after recovery with
\texttt{DETACH\_TIMEOUT} and \texttt{TERMINATION}; $T_{\text{detach}}$
itself has just restarted and cannot have elapsed. The expiry of an
approval, $T_{\text{approve}}$, is recorded as \texttt{valid\_until} in
\texttt{HUMAN\_RESOLVED}; recovery reads it back from the log, so an approval
that had lapsed before the restart stays lapsed.
\item \textbf{Recovery failure.} If the chain does not verify, the
bridge MUST NOT resume the session and leaves the key registry as it is.
The bridge quarantines the session (rule~Q): it
emits \texttt{RECOVERY\_FAILED} if it can still append, emits no
\texttt{TERMINATION}, and leaves the session to an operator.
A log that cannot be replayed at all (no \texttt{ATTACHED} event first, or
one without the fields the session is built from) quarantines its own
session in memory, appends nothing to it, and does not stop the recovery
of the other sessions. A chain
failure found during a re-attach is handled the same way, because the
cause is the same and is not attributable to the client.
\item \textbf{In-flight tool calls.} A call whose drain was committed
but whose delivery the runtime did not acknowledge before the crash
(\texttt{TOOL\_CALL\_DELIVERED}) is forwarded again, with the same
idempotency key (I13), when the session is back in
$(\textsc{Att},\textsc{None})$ (rule~D); it is not drained twice. A call
that was queued but not drained is subject to the usual drain checks. The
runtime deduplicates by key. Delivery is therefore at-least-once with
deduplication, never silently lost, and a call that was committed as
drained is never dropped silently either.
\end{enumerate}

Bridge availability (replication, failover to another host) is a
deployment concern beyond this rule: the rules above hold for any bridge
instance that has the durable log and the registered client keys.

\subsection{Additional invariants}\label{sec:additional-invariants}

The following invariants are conditioned on deployment-specific
concerns or strengthen the core set $\{I_1, \ldots, I_7\}$. They are
stated here to make normative requirements that would otherwise live
in \S\ref{sec:limitations} live instead in the specification.

\subsubsection{Event signing (Invariants I-SIGN-EVENT and I-SIGN-JWKS)}\label{sec:invariants-signing}

\begin{invariant}[I-SIGN-EVENT --- Event signing in GEC-less deployments (SHOULD)]
\label{inv:isign-sign}
In a GEC-less deployment (Definition~\ref{def:gecless}), the bridge
SHOULD sign each emitted event with an operator key whose JWKS
endpoint is published through an out-of-band channel bound to the
operator's deployment identity (per GAR \S 6.1).
In a GEC-present deployment, the bridge forwards unsigned events to
the GEC for signing; I-SIGN-EVENT is vacuous.
\end{invariant}

\begin{invariant}[I-SIGN-JWKS --- JWKS endpoint source restriction (MUST)]
\label{inv:isign-jwks}
The bridge MUST NOT accept a JWKS endpoint learned from
session-supplied metadata. The JWKS endpoint used by the bridge for
event signing (or for forwarding to the GEC) MUST be configured via
an out-of-band channel bound to the operator's deployment identity.
A JWKS endpoint learned from the client or from session-supplied
metadata is not a trust anchor.
\end{invariant}

I-SIGN-EVENT is SHOULD (not MUST) because event signing requires
operator key management infrastructure that may not be available in
testing or small deployments. Operators who omit signing accept
that downstream consumers cannot verify event provenance without an
out-of-band trust channel. I-SIGN-JWKS is MUST because accepting
session-supplied JWKS endpoints would allow the client (or an
adversary) to substitute a different signing key, breaking the
trust model; this is not optional regardless of deployment size.

\subsubsection{QUIC 0-RTT replay (Invariant I8)}\label{sec:invariants-quic}

\begin{invariant}[I8 --- QUIC 0-RTT replay rejection]\label{inv:I8}
The bridge MUST NOT accept a client re-attach whose underlying
transport is a replayed QUIC 0-RTT session. Replay detection is
governed by AGTP-SESSION \S 10.5, which forbids establishing a session in
0-RTT data, and by the replay-safety profile of AGTP-BINDINGS \S 4.6, which
restricts early data to safe idempotent methods \cite{I-D.hood-agtp-bindings};
the bridge defers to the transport layer for replay identification and rejects the re-attach (with
\texttt{RESUME\_REJECTED}) if the transport reports a replay.
\end{invariant}

I8 is conditioned on the transport layer's replay detection. In
deployments that do not use QUIC 0-RTT, I8 is vacuous. In deployments
that do, the bridge's responsibility is to \emph{honour} the
transport's replay verdict, not to perform replay detection itself.

\subsubsection{Challenge nonce lifetime (Invariant I9)}
\label{sec:invariants-ttl}

\begin{invariant}[I9 --- Challenge nonce lifetime]\label{inv:I9}
A challenge nonce (I11) is valid for $T_{\text{nonce}}$ after the bridge
issued it, measured on the bridge's monotonic clock (A6). A re-attach
that presents a nonce older than $T_{\text{nonce}}$ MUST be rejected
with \texttt{RESUME\_REJECTED} (reason \texttt{nonce\_expired}). The
rejection does not change the cell and does not terminate the session:
only a timeout ($T_{\text{detach}}$, $T_{\text{total}}$ or, while an
escalation is pending, $T_{\text{human}}$) ends a detached session. \emph{Default:}
$T_{\text{nonce}} = 60$~seconds. \emph{Configurable:} yes,
per-deployment.
\end{invariant}

A design that bounds the lifetime of a resume token issued at detach,
and terminates the session when it expires, cannot serve a client that vanished abruptly: the token is minted when
the client is already gone and has nowhere to be delivered. The
challenge is instead minted when the client asks for it, so its
lifetime can be short without ever stranding the client; I9 bounds the
window in which a captured nonce and signature could be replayed, and
nothing else.

\subsubsection{Authority scope versioning (Invariant I10)}
\label{sec:invariants-scope-version}

\begin{invariant}[I10 --- Authority scope versioning]\label{inv:I10}
The session's authority scope carries a monotonically increasing
\texttt{scope\_version} counter, incremented on each successful
narrowing (i.e., each \texttt{AUTHORITY\_NARROWED} event), in every
cell. The counter starts at 0 at attach. Each
\texttt{AUTHORITY\_NARROWED} event MUST include
\texttt{new\_version} = \texttt{previous\_version} + 1 and a scope that
is a strict subset of the current scope. A duplicate or out-of-order
narrowing, or one whose proposed scope equals the current scope, MUST be
rejected with \texttt{AUTHORITY\_NARROWING\_REJECTED}; a proposal that
widens the scope is governed by I3 and is rejected with
\texttt{AUTHORITY\_EXTENSION\_REJECTED}. A proposal that is not a strict
subset of the current scope is therefore rejected: if it contains any
permission the current scope lacks, including a partial overlap that also
drops permissions, it is an extension and is rejected as above. A scope change that happened
while the client was away is not a reason to refuse its re-attach: it is
part of the replay that I4 requires, and the client reconciles to the
new scope.
\end{invariant}

I10 gives auditors a monotonic history of authority scope changes per
session. Regulators can trace the scope trajectory from attach
(version 0) through any narrowings (versions 1, 2, \ldots) to the
current state. With I3 it makes the authority of a session
non-increasing for its entire life.

\subsubsection{Challenge-response resume (Invariant I11)}
\label{sec:invariants-token-auth}

\begin{invariant}[I11 --- Authenticated single-use challenge]\label{inv:I11}
\begin{enumerate}[leftmargin=2em,label=(\alph*)]
\item \emph{Key binding.} The client registers a public key at attach;
the key is recorded in the genesis \texttt{ATTACHED} event and bound to
\texttt{Owner-ID}.
\item \emph{Challenge.} To re-attach, the client first requests a
challenge, naming the key identifier it will sign with. The bridge
returns a fresh nonce of at least 128 bits from a
cryptographically secure source, bound to the session and to that key
identifier, valid for $T_{\text{nonce}}$ (I9) and usable once. The record of
the nonce holds the key identifier and whether the key was a registered,
unrevoked key of the session when the nonce was issued. The answer to a
challenge is the same for every key identifier, registered or not: a nonce of
the same form and lifetime, counted against the same limit (f); no answer to
a challenge depends on the registry. A nonce issued for a key that was not
registered at issue never verifies, however the registry changes later, and a
nonce presented with another key identifier than the one it was issued for
does not verify either; both are answered as a request whose signature does
not verify (\texttt{bad\_signature}), recorded as \texttt{nonce\_binding} and
not spent. Issuing a challenge is not a
governance event and changes no state. The bridge SHOULD return a uniform rejection
to unauthenticated callers (e.g., for an unknown session ID) so
that the \texttt{challenge} op does not confirm session existence
or lifecycle state. Nonces are not durable (Section~\ref{sec:recovery}).
\item \emph{Proof of possession.} The re-attach request carries
$(\mathit{nonce}, \mathit{seq}, \mathit{hash})$ and a signature, under
a registered client key, over the tuple of Appendix~\ref{app:wire-format}
(signature scope), which binds the operation (re-attach or join), the
session, the key id, the nonce and the prefix. The request is authentic
iff the signature verifies for the operation received. A request that
names a key id the session does not register is answered as one whose
signature does not verify (\texttt{bad\_signature}), so the answer does not
list the registered key ids; the log records \texttt{unknown\_key}. A
registration of a further key (\texttt{add\_client}) is verified in the same
order: the signature first, and the lifecycle of the session (unknown,
terminated, quarantined, key-exhausted) is named only to a caller whose
signature verified; any other caller is answered \texttt{bad\_signature}.
\item \emph{Single use.} A nonce is consumed by the first authentic
request that presents it, whether that request is then accepted or
rejected for a later reason; a consumed nonce is never accepted again.
A request that is not authentic does not consume the nonce. A request
refused with \texttt{resume\_in\_progress}, \texttt{already\_attached} (a
re-attach on an attached session) or \texttt{not\_attached} (a join on a
detached one) is authentic and so spends its nonce: a second device that
meets a re-attach in progress asks for a new challenge. These reasons are
given only after the key, the signature, the revocation, the binding of the
nonce and the nonce itself have passed, so none of them tells an
unauthenticated caller anything about the attachment. A request that was
presented and refused is never accepted later, whoever presents it again and
whatever the registry has become in between.
\item \emph{No state-changing rejection.} A request that fails (c), (d),
I9, I8, I11-REVOKE or I4 MUST be answered with \texttt{RESUME\_REJECTED} and MUST NOT
change the cell. Through the re-attach path, only a failure of hash-chain
verification after acceptance (\texttt{RESUME\_FAILED}), which quarantines
the session (rule~Q), or a timeout, which terminates it, changes the
session's standing.
\item \emph{Bounded issue, no eviction.} The bridge keeps at most a
fixed number of live nonces per session (the reference bridges use 16), where
live means unspent and unexpired. Issuing a challenge MUST NOT
invalidate a live nonce: when the limit is reached the bridge discards
spent and expired records to make room and, if none exists, refuses the
challenge. A nonce a genuine client already holds therefore stays valid
whatever other callers request. The \texttt{challenge} op is
unauthenticated, so the limit also bounds what a caller without a key can
do; the refusal reveals only that the session cannot issue a nonce now,
and the rejection of a re-attach request is unchanged (the uniform rejection
of (b)).
\item \emph{Bounded logging of unauthenticated failures.} A failure
before authentication (\texttt{transport\_replay}, \texttt{unknown\_key},
\texttt{nonce\_binding}, \texttt{bad\_signature}) is recorded as \texttt{RESUME\_REJECTED} only
until the session has recorded a fixed number of such failures since its
last authentic request (the reference bridges use 16); further ones are
refused in the same way and not recorded. A request whose signature
verifies restores the allowance, and every failure of such a request is
recorded. A replay of a captured authentic request also verifies, so
whoever holds its bytes, not only the holder of the key, can restore the
allowance: the limit bounds the failures before authentication between two
verifying requests, not the log as a whole (Residual denial, below). The limit is held in
memory, so a restart restores it; the number of unauthenticated records
a session can gain is bounded by the limit for each period between
restarts and authentic requests.
\end{enumerate}
\end{invariant}

The verifier checks in this order: (1) key and signature, then
(2) revocation (I11-REVOKE), then (3) the binding of the nonce to the key
and to the registry at issue, then (4) nonce (unknown, consumed,
expired), then (5) prefix (I4), and only then what depends on the state of
the session (an attachment in progress, an attached or a detached session),
so that an
unauthenticated party learns nothing about the state of nonces, of the
chain or of the attachment beyond a rejection. A timer that expires in the
very tick of a request waits for the request; one that expired earlier has
ended the session, and the request is refused uniformly. I11 ensures that knowing the session
identifier is not enough and that a captured request cannot be replayed:
one that was accepted has spent its nonce, an unspent nonce expires after
$T_{\text{nonce}}$, and a request refused before authentication cannot
become acceptable by a later registration of its key, because a nonce
issued for a key that was not registered never verifies. A rejected request does not change the cell, so a
party without a key cannot terminate a session by presenting bad requests;
it can still end a detached session by holding the live-nonce limit
(Residual denial, below). Theft of a client key is out of scope (A7).

\paragraph{Residual denial.} Item (f) stops a caller without a key from
invalidating a nonce, and item (g) bounds the failures before
authentication that it can have recorded between two verifying requests.
Neither bounds the log as a whole: a caller that holds the bytes of one
captured authentic request can replay it, each replay verifies, is
recorded (\texttt{nonce\_used}) and restores the allowance of (g), so the
log grows linearly in the caller's effort (Appendix~\ref{app:attacker},
Limits). Nor do they stop it from using up the live-nonce limit: while the limit is
held by its nonces, a client that needs a new challenge is refused until
one expires, up to $T_{\text{nonce}}$ later, and the caller must keep
asking at that rate to keep the limit held. A client that already holds a
nonce is not affected. A detached client that needs a challenge is, and if
the caller keeps the limit held until $T_{\text{detach}}$ elapses, the
session ends with \texttt{DETACH\_TIMEOUT}: a party that knows only the
session identifier can terminate a detached session this way. L1 is
therefore a bound on how long a session stays detached, not an
availability property. A deployment that needs more than this authenticates
or rate-limits the transport in front of \texttt{challenge} (for example
per source address); the contract has no identity to limit by before a
signature.

\paragraph{Outsider.} In this document an \emph{outsider} is a party that
observes and modifies the traffic between a client and the bridge, knows
the session identifier, and has seen earlier challenges and requests, but
holds no registered private key of the session (A7). I11 states that an
outsider cannot have a re-attach request accepted, and cannot spend a
nonce, except with a probability that is negligible under the assumptions
of A3 on the signature scheme and on the nonce source; the reference
bridges and the model meet this also for a request that was refused
before its key was registered (above). I11-REVOKE extends
the refusal to a party that holds a key revoked before the request: its
signature verifies, and the request is refused as \texttt{revoked\_key}.
The model has an explicit attacker (Appendix~\ref{app:attacker}): it observes
the log and the wire, replays and reorders, submits with keys it does not
hold and never obtains a client key. TLC checks, in the attacker
configuration, that every request accepted on entering $a=\textsc{Res}$ or
on a join carries a signature that an honest client produced for that key,
nonce and prefix verdict (\texttt{NoForgery}), so the attacker cannot build
an acceptable signature of its own, and that every accepted request was
presented by an honest client and that a request that was refused is never
accepted later (\texttt{SP\_Honest}, \texttt{SP\_NoReplay}). The unforgeability of the signature scheme is assumed
(UF, stated there) and the statement about concrete signatures rests on the
tests and mutants of the reference implementations.

\paragraph{Relayed requests.} The signed message of a re-attach covers
the operation, the session, the key, the nonce and the prefix, and no
identifier of the transport channel. A party that relays a genuine client's signature,
obtained for a nonce the bridge issued, is therefore not distinguishable
from the client at the contract level. Protection rests on the client
signing only nonces it requested from the bridge it means to reach, and
on the transport authenticating the bridge to the client. A deployment
that needs more can bind the signature to a channel identifier (for
example a TLS exporter value); this document does not specify that
binding.

\paragraph{Extension I11-REVOKE (key lifecycle).}
A key that has been revoked (recorded in the log with
\texttt{KEY\_REVOKED}) MUST NOT be accepted for re-attach or join,
even if the signature verifies. The rejection is
\texttt{RESUME\_REJECTED} (reason \texttt{revoked\_key}); the cell
does not change, and the nonce is spent (I11d). The verifier checks
revocation after key/signature verification and before the nonce
check, so an unauthenticated party learns nothing about nonce state.
Spending the nonce is a deliberate trade: leaving it live on a
revoked-key attempt would tell the holder of a revoked key that the
nonce is still outstanding (an oracle), at the price that such a holder,
who has a valid signature, can burn the outstanding nonce and force the
legitimate client to request another; this costs one round trip and
grants no access.

Revocation is operator-signed (Section~\ref{sec:invariants-token-auth}): the \texttt{KEY\_REVOKED}
event carries an Ed25519 signature under the operator key over
\texttt{["bc/v1/revoke", bridge\_id, sid, key\_id, reason, seq, head]},
and the bridge MUST verify it before committing the event (reason
\texttt{operator\_signature\_invalid}; nothing is committed). A client
may \emph{request} revocation (\texttt{KEY\_REVOCATION\_REQUESTED},
informational), but only a revocation with a verifying operator signature
is executed. The same holds for \texttt{OPERATOR\_ADD\_KEY}
(\texttt{["bc/v1/add-key", bridge\_id, sid, key\_id, scheme, public\_key,
seq, head]}) and \texttt{OPERATOR\_TERMINATION}
(\texttt{["bc/v1/terminate", bridge\_id, sid, reason, seq, head]}). The
bridge identity and the log position bind a signature to one bridge and
one moment: \texttt{seq} and \texttt{head} are the sequence number and the
hash of the last committed event of the session when the operator signs,
and the bridge verifies against the current last event, so a signature
verifies only while nothing has been committed after the position it names.
A captured operator action cannot be replayed at another bridge or later in
the session's life; if the log has moved, the operator signs again. The
bridge checks them against the operator key that
verifies scope grants; which key that is, and how it is protected, is
the deployment's.

\textbf{Soft vs hard revocation.} The \texttt{reason} field
distinguishes the two:
\begin{description}[leftmargin=2em]
\item[\texttt{rotation} (soft).] The key is revoked for new
re-attach and join, but already-attached connections using this key are
not disconnected and the key stays in $C(s)$ until its transport leaves
(I15(e)); if the key is the one whose re-attach is being verified, the
verification is abandoned as for a hard revocation. Used for rotation (\texttt{CLIENT\_ADDED} $\to$
overlap $\to$ \texttt{KEY\_REVOKED}).
\item[\texttt{compromise} (hard).] The key is revoked for all
purposes. A connection using this key leaves $C(s)$ at once (I15(e)); if
it was the only connected client, or the key whose re-attach was being
verified (a soft revocation of that key does the same), the session is forcibly detached (\texttt{DETACHMENT} +
\texttt{RESUME\_PREPARED}). Used for compromise response.
\end{description}

\textbf{Key exhaustion.} When the last active (unrevoked) key in $K(s)$ is
revoked, the session enters a \texttt{key\_exhausted} state: a side
predicate on the session, not a cell of the product machine and not
rule~Q. The session keeps its current cell and accepts no client
operation. The timers keep running: $T_{\text{detach}}$,
$T_{\text{human}}$ and the cumulative cap $T_{\text{total}}$ are
evaluated as in any session, a restart treats the session like any other
(it is recovered as detached and $T_{\text{detach}}$ restarts), and the
loss of a client's transport detaches it. The rule of permitted operations
governs the operations that a client or the operator requests; the events the
timers drive still fire. Only these operations are permitted: \texttt{OPERATOR\_ADD\_KEY} (adds a new key, clears
\texttt{key\_exhausted}) and \texttt{OPERATOR\_TERMINATION}
(transitions to $\bot$). A termination without the operator's signature
(the unsigned \texttt{terminate} of the transport binding) is refused on an
exhausted session and commits nothing; a quarantined session is still ended
by it. The operator's signed termination of a session that is both
quarantined and key-exhausted appends \texttt{OPERATOR\_TERMINATION}, then
\texttt{TERMINATION}, as for any exhausted session. Unless the operator adds a key, an exhausted
session that is detached therefore terminates after the detachment bound
($T_{\text{detach}}$ of the detachment, or $T_{\text{total}}$); one that
stays in $a=\textsc{Att}$ on a softly revoked key does not (see the hazard
below). A hard revocation of the last connected key also detaches the
session, as defined in I11-REVOKE. Once
the session has terminated, adding a key or terminating it changes
nothing. This is distinct from rule~Q (storage
corruption): key exhaustion is a deliberate action by an authorised
party, and it has a recovery path.

\textbf{Deployment hazard: rotating the only key.} A soft revocation
of the only registered key reaches \texttt{key\_exhausted} just as a
hard one does, but leaves the revoked key connected: the session stays
in $a=\textsc{Att}$, accepts no client operation, and runs no detach
timer. It ends only when the transport of that client is lost (the
session is then detached and $T_{\text{detach}}$ applies) or the operator
acts, so the runtime behind it may keep executing meanwhile. The
recommended operator procedure is to rotate by adding the new key first
(\texttt{CLIENT\_ADDED}, then \texttt{KEY\_REVOKED} of the old key), so
that the registry never holds zero active keys. A hard revocation of
the only key detaches the client, so the exhausted session is detached
and ends by timeout unless a key is added in time.

\subsubsection{Durable append (Invariant I12)}
\label{sec:invariants-durable}

\begin{invariant}[I12 --- Durable append]\label{inv:I12}
An event is \emph{committed} when it is durable on the medium assumed
by A1 (e.g., after \texttt{fsync} or an equivalent barrier). The bridge
MUST NOT change the cell, forward a message to the agent runtime, or
acknowledge any client request that depends on an event until that event
is committed.
\end{invariant}

I12 is what gives I7 its meaning in a process that can crash. It is
checked by the recovery rules of Section~\ref{sec:recovery}.

\subsubsection{Queue integrity and approval freshness (Invariant I13)}
\label{sec:invariants-queue}

\begin{invariant}[I13 --- Queue integrity and approval freshness]\label{inv:I13}
Every queued tool call carries an \emph{idempotency key}, unique within
the session, which MUST be forwarded to the runtime with the call. The
key is supplied by the runtime when it can and generated by the bridge
otherwise; a call that is forwarded again after a restart reuses the
key it was queued with, never a new one. The
queue has a maximum size ($Q_{\max}$, default 64) and each queued call a
lifetime ($T_{\text{call}}$, default $T_{\text{human}}$, with
$T_{\text{call}} \leq T_{\text{human}}$ as a
MUST, Section~\ref{sec:timeouts}). A call arriving at a full
queue MUST be rejected with \texttt{TOOL\_CALL\_REJECTED}; a call whose
key is already queued MUST NOT be queued twice; an expired call MUST be
removed with \texttt{TOOL\_CALL\_EXPIRED}. A human approval
(\texttt{HUMAN\_RESOLVED}) is valid for $T_{\text{approve}}$ (default
15 minutes) from its emission, strictly before it has elapsed; the expiry is
recorded for each call queued under the escalation when the approval is
given, and a later approval does not extend it; a call held under the detached-call
policy (rule~D, calls outside an escalation) has no escalation approval
and relies on its own answer (I14) and $T_{\text{call}}$. When the queue drains (rule~D), every
call MUST be re-validated against the \emph{current} authority scope and
approval, and against its answer (I14); a call that requires a permission the scope no longer
contains, or whose approval has expired, MUST NOT be forwarded and is
audited with \texttt{TOOL\_CALL\_REJECTED} or \texttt{TOOL\_CALL\_EXPIRED}
respectively.
\end{invariant}

I13 prevents a long detach from replaying stale or over-privileged side
effects (for example, a queued \texttt{git push} that was approved
hours earlier). A rejected-at-drain call is not retried; the agent must
re-escalate.

\subsubsection{Answer binding (Invariant I14)}
\label{sec:invariants-answer}

\begin{invariant}[I14 --- Answer binding]\label{inv:I14}\label{inv:answer_binding}
Let $h(x)$ be the SHA-256 of the canonical JSON of $x$. Every queued tool
call records $h_q = h(\{\mathit{tool}, \mathit{permission},
\mathit{content}\})$ in \texttt{TOOL\_CALL\_QUEUED}, where \emph{content}
is what the human is shown to decide on that call (arguments and any
explanation), taken as null when the runtime supplies none; the tool name
and the permission are hashed in every case, so the hash does not depend on
what the runtime chose to put in \emph{content}. The event retains the
tool, the permission and the content as hashed. A human decision on a call is
recorded as \texttt{TOOL\_CALL\_ANSWERED} carrying the call's key, the
decision (approve or deny) and $h_a$, the hash of the call (tool, permission
and content) the human's client displayed when the human decided. The first answer to a
call is final; a second one MUST be refused. When the queue drains
(rule~D), a call MUST be forwarded only if it has an approving answer
with $h_a = h_q$, and the request the bridge forwards MUST be the tool,
permission and content the bridge hashed at queue time, taken from the
retained record (the bridge MUST NOT synthesise, substitute or alter the
call between queue and forward). Otherwise it MUST NOT be forwarded and is audited with
\texttt{TOOL\_CALL\_REJECTED} and the reason \texttt{unanswered},
\texttt{denied} or \texttt{answer\_mismatch}. The
\texttt{TOOL\_CALL\_DRAINED} event repeats $h_q$ together with the
retained tool, permission and content, so a forward repeated after a restart
sends the same request.
\end{invariant}

I14 binds the bridge's forwarded request to the hash recorded at
queue time and the hash reported in the answer. A view
that is reset, re-rendered or compacted after a resume cannot silently
change what was approved: if the client now displays something else, its
answer carries a different hash and the call is not forwarded. I14
protects against content drift by an honest client whose rendering
differs from what the bridge recorded; it does not protect against a
malicious client that reports whatever hash it was given (the bridge has
no independent view of what the client displayed). The
approval of I13 (\texttt{HUMAN\_RESOLVED}) says that the escalation is
over and how long answers stay fresh; I14 says what each answer was an
answer to. The two anchors are independent and a drain needs both: a call
is forwarded only if its own answer matches (I14) and the approval of its
escalation has not lapsed (I13). With the persistent profile an approval
given while the client is away is void $T_{\text{approve}}$ after the
human gave it, so a client that returns later finds the call expired
(\texttt{approval\_ttl}) and the agent must re-escalate; this is
deliberate. I14 does not say that the content is a faithful description
of the action the runtime will perform: that is outside what a bridge can
see (Section~\ref{sec:coverage}).

\subsubsection{Client set (Invariant I15)}
\label{sec:invariants-client-set}

\begin{invariant}[I15 --- Client set coherence]\label{inv:I15}
\begin{enumerate}[leftmargin=2em,label=(\alph*)]
\item \emph{Derivation.} $C(s) \subseteq K(s)$, and the attachment is
the function of $C(s)$ defined in Section~\ref{sec:product-model}:
$a=\textsc{Att}$ iff $|C|\geq 1$ (outside a verification in progress).
\item \emph{Departure.} The transport disconnect of a client while
$|C| \geq 2$ MUST emit \texttt{CLIENT\_LEFT} and MUST NOT change the cell,
emit \texttt{DETACHMENT} or start $T_{\text{detach}}$. Only the
departure of the last connected client detaches the session (I1).
\item \emph{Joining.} After attach (whose key is the first member of $C$), a key enters $C$ only through an authentic request
that passes I9, I11 and I4 (a join while $a=\textsc{Att}$, a re-attach
while $a=\textsc{Det}$). A re-attach request while $a=\textsc{Att}$ is
refused with \texttt{already\_attached}, so a proof made for a re-attach
never causes a join. A key not in $K$ is never accepted.
\item \emph{Registration.} A key enters $K$ only at attach, through
\texttt{CLIENT\_ADDED} or through the operator-signed
\texttt{OPERATOR\_ADD\_KEY} (on a \texttt{key\_exhausted} session, in any
cell). \texttt{CLIENT\_ADDED} is authorised by a client in $C$ by a signature
under its own registered key over the new key, the session and a
bridge-issued single-use nonce, under its own domain tag (Appendix~\ref{app:wire-format});
a request without such a signature, with one made for another operation,
key or session, or with a spent, unknown or expired nonce is refused with
a named reason and commits nothing. Knowing the key id of a connected
client is therefore not authority. \texttt{CLIENT\_ADDED} is impossible while
$a \neq \textsc{Att}$. The new key is bound to the same
\texttt{Owner-ID}.
\item \emph{Revocation (I11-REVOKE).} Revoked keys remain in $K(s)$
with \texttt{revoked\_at} and \texttt{reason} markers; they are
never deleted. A key revoked for \texttt{compromise} (hard) leaves
$C(s)$ at once. A key revoked for \texttt{rotation} (soft) stays in
$C(s)$ for as long as its transport is connected and leaves it when the
transport does; it cannot join or re-attach. A key enters the revoked
state via \texttt{KEY\_REVOKED} (operator-signed, see the
revocation paragraph of Section~\ref{sec:invariants-token-auth}) and
never returns to the active state.
\end{enumerate}
\end{invariant}

I15 is what makes ``attached'' mean the same thing for a phone and a
laptop: the session is detached only when nobody is connected, a scope
narrowed from one device reaches every other device through the replay of
its next join or re-attach (I4), and no device can be added by someone who
is not already trusted by an attached one. With I11-REVOKE, a hard-revoked key is no longer in $C(s)$ even if its
transport is still connected (hard revocation removes it, and detaches the
session if it was the last connected client), whereas a
soft-revoked key stays in $C(s)$ until its transport leaves; the session
is detached only when no key in $C(s)$ has a connected transport.

\subsubsection{Checkpoint anchoring (Invariant I-CHECKPOINT)}
\label{sec:invariants-checkpoint}

\begin{invariant}[I-CHECKPOINT --- Checkpoint anchoring (SHOULD)]\label{inv:icheckpoint}
The bridge SHOULD emit a \texttt{CHECKPOINT} event at least every $N$
events (default 1000) or every $T_{\text{cp}}$ (default 5 minutes),
whichever comes first. The event carries the Merkle tree hash of RFC~9162, \S2.1.1
(Appendix~\ref{app:merkle}) over the hashes of all preceding events of the
session, the tree size, and
a signature under the operator key. The bridge SHOULD submit the
checkpoint to an external append-only transparency log (e.g., SCITT)
and record the receipt in the event payload.
\end{invariant}

Checkpoints let an auditor detect a bridge that rewrites history after
the fact, which the hash chain alone cannot (A4), provided the receipt is
anchored in an external log (a SHOULD); without it the operator key can
sign a rewritten log afresh. They also let recovery and
re-attach verify only the chain suffix after the latest verified
checkpoint, giving verification cost $O(m)$ in the suffix length $m$
rather than $O(n)$ in the whole log. The $O(n)$ full-chain verification
remains the fallback.

\subsection{Product model and derived view}\label{sec:mece-status}

The product machine of Definition~\ref{def:sm} is mutually exclusive and
collectively exhaustive over the dimensions it names: every
configuration of a live session is exactly one of the active cells
or $\bot$, and the human and attachment dimensions are independent. In
particular there is no artificial transition that merely changes which
dimension is salient. The names of the derived view are
retained because they make the gap classes G1--G4 of
Definition~\ref{def:gap} visible at a glance, but they are a lossy fold
and are not normative. Appendix~\ref{app:product-model} gives the
mapping.

\subsection{Safety argument}\label{sec:safety-argument}

The following is a \emph{safety argument}, not a formal theorem. A
formal theorem would require a proof (by induction, model checking, or
both) that $\forall \tau: \mathrm{Safe}(\tau)$ holds under the
invariants and axioms, for the predicate $\mathrm{Safe}(\tau)$ over
execution traces $\tau$ defined below. The
present argument is informal and goes property by property, not by a
full induction over the trace length;
the conformance suite in Section~\ref{sec:conformance} provides
executable evidence for the invariants in a controlled test setting,
but does not constitute a formal proof.

Let $\tau$ be an execution trace of a bridge $B$. Define
$\mathrm{Safe}(\tau) = \mathrm{Audited}(\tau) \wedge
\mathrm{NoDrift}(\tau) \wedge \mathrm{NoLoss}(\tau)$, one conjunct for
each of the conclusions (a), (b) and (c) of Safety Argument~1:
\begin{description}[leftmargin=2em]
\item[$\mathrm{Audited}(\tau)$.] Every step of $\tau$ that changes the
cell also commits events whose \texttt{prev\_hash} links them to the
previous committed event (I5, I7, I12).
\item[$\mathrm{NoDrift}(\tau)$.] No step of $\tau$ widens the scope
that the bridge holds (I3, I10).
\item[$\mathrm{NoLoss}(\tau)$.] Every event committed before a detach, a
restart or a re-attach is replayed to the re-attaching presenter or is
recoverable from the committed log (I1, I4, I5, I11), for a log whose
last events the bridge has not truncated or rewritten (A1, A4; I5 does not
detect that).
\end{description}
TLC checks these conjuncts on bounded configurations, and not all
equally. $\mathrm{NoDrift}$ is checked directly: I3 and I10 are TLC state
invariants and step properties. In $\mathrm{Audited}$ the step property
of I7 (a cell change commits new events) is checked, while the hash
link is a corruption flag in the model (I5). In $\mathrm{NoLoss}$, I1 is a
TLC state invariant, while the prefix check is a flag (I4) and the proof of
possession is a possession predicate against an attacker that holds no client
key (I11). Table~\ref{tab:invariant-independence} records the
status of each invariant.

\begin{quote}
\textbf{Safety Argument 1} (no unaudited state transition, no
authority drift, no lost events). Let $B$ be a bridge implementation
that satisfies Invariants I1--I7, I10, I11 and I12 under Axioms A1--A7. Then
no execution of $B$ produces (a) an unaudited state transition,
(b) an authority extension granted in any cell, or (c) a lost event during
detach--resume--reattach or across a restart, for a committed log whose
last events have not been truncated or rewritten (A1, A4).
\end{quote}

\emph{Argument.} The argument takes each conclusion from the invariants
and lemmas it names; it is not a complete induction over the trace length.

\emph{Base case.} The genesis event is emitted with
\texttt{prev\_hash} = $\epsilon$ (empty string), satisfying I5, and is
committed before the cell is set to $(\textsc{Att},\textsc{None})$,
satisfying I7 and I12. No authority extension occurs at attach, satisfying
I3. No detach has occurred, so I1 is vacuously satisfied. The tool queue
is empty, so I2 is vacuously satisfied.

\emph{Per-property argument.} Assume the invariants hold. It does not
split into the kinds of step (timeout termination, drain and forward,
operator actions, client departure, quarantine by rule~Q), and those step
kinds are not covered by it; they rest on the invariants' own statements, on
the tests and mutants, and on TLC on bounded configurations.

\emph{(a) Unaudited state transition.} By Lemma~\ref{lem:atomicity}
(I5 $+$ I7 $+$ I12), the events are committed before the cell changes
and carry a valid \texttt{prev\_hash} linking them to step $k$'s last
event. Hence the transition is audited.

\emph{(b) Authority extension.} By Lemma~\ref{lem:no-drift}
(I3 $+$ I10), a widening attempt in any cell is rejected, emits
\texttt{AUTHORITY\_EXTENSION\_REJECTED}, and does not mutate the scope
the bridge holds. Hence no authority extension is granted; one that the
runtime performs on its own is reported, not prevented.

\emph{(c) Lost event during detach--resume--reattach.} By
Lemma~\ref{lem:recover} (I1 $+$ I4 $+$ I5 $+$ I11), on detach,
\texttt{RESUME\_PREPARED} is committed before the cell changes. On
re-attach, only a presenter that holds a registered key and answers a
live, unused nonce is admitted, and only for a history that is a prefix
of the chain (I4, I11); every later event is replayed to it. On success,
the bridge enters $a=\textsc{Res}$ and re-verifies the hash chain. By I5, any edit or reordering of the chain before its last events is
detected; removal or rewriting of the last events is not (Section~\ref{sec:limitations}). Across a restart, the cell is rebuilt from the committed log
(Section~\ref{sec:recovery}), so no committed event is lost, provided
the committed log has not been cut back or had its last events rewritten
(A1, A4). $\square$

\begin{remark}
The argument above is informal. The invariants it relies on are checked
by TLC on bounded configurations (Appendix~\ref{app:formal}) to the
extent recorded in Table~\ref{tab:invariant-independence}: some are
checked as TLC state invariants or step properties, some are abstracted
in the model (I4, I5 and I14 hold only as flag abstractions there; I11 is
checked against an attacker, with the scheme's unforgeability assumed);
I11-REVOKE is checked in part (the state invariant \texttt{IRev} and
step properties; the operator signature and the wire form only by tests),
and I8 is tests-only. For I6 and for the no-eviction rule of I11 (f) the
TLC checks are vacuous as configured, so these rest on the tests and
mutants. The reference suite, its
mutants and the shared vectors (Section~\ref{sec:conformance})
exercise the invariants in a controlled test setting. The invariants
are the author's statement of what the bridge must preserve, derived from
the design; the STRIDE-lite threat model (Section~\ref{sec:threat-model}) is
a second enumeration by the same author, and no hazard analysis from the
operator's side backs the set. Consistency with the stated invariants
therefore does not show that the set is complete for an operator's safety.
A full formal
proof requires a proof of $\mathrm{Safe}(\tau)$ for every trace of
the unbounded product machine, and Coq/Isabelle mechanisation;
these are targets for a future revision.
\end{remark}

% ================================================================================
